% ============================================================
%   Chapter 7  --  Reducing Shield Dependence
% ============================================================

\chapter{Reducing Shield Dependence}\label{ch:dependence}

The evaluation protocol of Chapter~\ref{ch:experimentation} makes an explicit concession: the
shield is kept active at evaluation time because ``a policy trained with a shield has learned to
rely on its interventions''. This chapter takes that concession as its research question. It
first \emph{measures} how dependent the trained agent actually is
(Section~\ref{sec:dep-problem}--\ref{sec:dep-probe}), then studies two strategies to reduce that
dependence: a Large-Language-Model (LLM) meta-controller that dynamically loosens the shield
(Section~\ref{sec:dep-llm}), and a behavioural-cloning scheme that turns the shield into a
\emph{teacher} (Section~\ref{sec:dep-teacher}). The first strategy is a carefully evaluated
\emph{negative} result; the second succeeds to the point where the agent drives \emph{better}
without the shield than with it. Section~\ref{sec:dep-collapse} analyses an instructive failure
of the teacher scheme, and Section~\ref{sec:dep-discussion} distils the central lesson: shield
dependence is a \emph{learning} problem, not a \emph{scheduling} problem.

\section{The shield-dependence problem}\label{sec:dep-problem}

Two design choices, each reasonable in isolation, conspire to make dependence the optimal
behaviour for the agent.

\paragraph{No cost for being shielded (R1).} The reward shaper of
Section~\ref{sec:design-reward} contains no term that penalises shield activation. When the
shield overrides a reckless action it keeps the vehicle on the road and moving, so the agent
collects \emph{more} return than if its own unsafe action had reached the simulator. Being
shielded is therefore free and, in fact, beneficial: the reward landscape makes the dependent
policy reward-optimal.

\paragraph{No gradient where it matters (R2).} The masked PPO objective
(Section~\ref{sec:design-agent-ppo}) computes the policy gradient only over unshielded steps, weighting
each transition by \((1-\alpha_t)^2\), where \(\alpha_t\in[0,1]\) is the shield intensity. The
shielded steps are precisely the dangerous, near-boundary states. The actor therefore receives a
learning signal only on the easy states the shield never touched, and is \emph{structurally
prevented} from learning how to act at the boundary. The critic does observe those states (its
loss is unmasked), but the actor cannot act on them. The framework even stores the shield's
executed action and recomputes its log-probability, ideal imitation data, and then the
mask discards it.

Together, R1 and R2 produce an agent that ``rides the shield boundary'': it calibrates its own
imprecision to whatever tolerance the shield currently allows, so its competence \emph{with} the
shield says nothing about its competence \emph{without} it.

\section{Measuring dependence: the shield-off probe}\label{sec:dep-probe}

Quantifying dependence requires an evaluation that the shaped reward and the shield-activation
rate both hide. Two instruments are added.

\paragraph{Shield-off probe.} Every \(N_{\mathrm{probe}}\) training episodes the shield is placed
in a \emph{bypass} mode in which it forwards the proposed action to the simulator unchanged, and
the deterministic policy is rolled out for a handful of episodes with no buffer writes and no
optimiser steps; the observation normaliser is frozen and restored afterwards so the probe does
not contaminate training. The resulting shield-off success, crash and off-road rates are the
\emph{true} learning curve: they answer ``can the agent drive on its own yet?''.

\paragraph{Dynamic\,/\,static reliance split.} The per-episode intervention count is separated
into a \emph{dynamic} rate (vehicles and pedestrians) and a \emph{static} rate (barriers, walls,
guardrails). On the Town04 highway the static channel dominates by roughly \(30{:}1\), so the
overall shield-activation rate is mostly a function of the map geometry rather than of the
agent's skill; the split exposes the part of reliance that is actually learnable.

\section{First approach: an LLM meta-controller}\label{sec:dep-llm}

\subsection{Design}

The first hypothesis, inspired by the use of LLMs as meta-optimisers for reward
design~\cite{ma2023eureka} and by automatic curriculum learning, is that an LLM can
\emph{wean} the agent by progressively loosening the shield as the agent improves. A single
scalar \emph{permissiveness} \(p\in[0,1]\) controls every tunable threshold \(\tau_i\) of the
adaptive shield by linearly interpolating each one between a strict and a permissive endpoint, so
that \(p=1\) recovers the production shield and \(p=0\) the maximally weaned shield
(Figure~\ref{fig:permissiveness-ramp}; the interpolation is Equation~\eqref{eq:permissiveness} in
Appendix~\ref{app:math}).

\begin{figure}[H]
  \centering
  \includegraphics[width=0.8\textwidth]{imgs/permissiveness_ramp}
  \caption{The permissiveness scalar \(p\) sets every tunable shield threshold \(\tau_i\) at once by
           linear interpolation between its strict value (at \(p=1\), the production shield) and its
           permissive value (at \(p=0\), maximally weaned). Three representative thresholds are
           shown; the LLM meta-controller proposes a single \(p\) every \(100\) episodes and the
           deterministic safety governor accepts or vetoes it.}
  \label{fig:permissiveness-ramp}
\end{figure}

Every \(100\) episodes, a local LLM (Gemma~\cite{team2024gemma} via Ollama, running on CPU so as not to contend with
CARLA for GPU memory) receives a compact snapshot of the safety and competence metrics and
proposes a target permissiveness. Crucially, the LLM never acts directly: a deterministic
\emph{safety governor} mediates every proposal under four hard rules: a safety floor that
forces the shield \emph{stricter} whenever crash or off-road rates breach a ceiling, a ratchet
that only loosens when the metrics are below stricter targets, a bounded per-decision step, and a
dwell time between changes. The governor makes ``safety first'' a property of code rather than of
the language model.

\subsection{Result: loosening is not weaning}

A \(3{,}500\)-episode run with this controller produced three independent and mutually
reinforcing observations.

\begin{enumerate}
  \item \textbf{The LLM added no adaptive intelligence.} It proposed ``loosen'' on all \(34\)
        decisions, irrespective of the crash rate; the deterministic governor performed all of
        the safety regulation (twelve \emph{hold} and two \emph{force-tighten} vetoes). A
        three-line fixed schedule would have behaved identically.
  \item \textbf{Loosening did not reduce reliance.} As Figure~\ref{fig:llm-negative} shows, the
        permissiveness fell from \(1.0\) to \(\approx0.55\) while the shield-activation rate
        remained essentially flat: the Pearson correlation between the two series is
        \(\approx 0\). Meanwhile the correlation between permissiveness and the crash rate is
        \(-0.50\): removing the shield's help immediately surfaced unsafe behaviour, exactly the
        ``boundary-riding'' dynamic of Section~\ref{sec:dep-problem}.
  \item \textbf{Reliance was geometry, not skill.} \(96.7\%\) of all interventions were static
        (guardrail proximity on the Town04 highway), a structural feature of the map that the
        tunable thresholds barely affect.
\end{enumerate}

\begin{figure}[H]
  \centering
  \includegraphics[width=0.85\textwidth]{imgs/llm_negative}
  \caption{LLM meta-controller over \(3{,}500\) episodes. The controller loosens the shield
           (permissiveness, left axis) but the shield-activation rate (right axis) does not fall:
           the agent compensates by driving less precisely, keeping its reliance constant. The
           crash rate, in contrast, rises as the shield is loosened.}
  \label{fig:llm-negative}
\end{figure}

The conclusion is unambiguous: loosening the shield's \emph{thresholds} does not teach the agent
anything; it merely moves the boundary the agent rides into, paying the safety cost of weaning
without buying any independence. The LLM, additionally, never justified its place over a
deterministic schedule.

\section{Second approach: the shield as a teacher}\label{sec:dep-teacher}

\subsection{Filling the gradient hole}

If dependence is caused by the absence of gradient on the shielded states (R2), the remedy is to
supply one. The masked PPO surrogate is left untouched, since it is \emph{correct} not to apply an
importance-weighted update to an off-policy action that the shield substituted, and a
supervised behavioural-cloning (BC) term is \emph{added} on exactly those masked steps, pulling
the policy mean towards the action the shield executed:
\begin{equation}\label{eq:bc}
  \mathcal{L}_{\mathrm{BC}}(\theta) \;=\;
  \frac{\sum_t \alpha_t\,\bigl(\tanh(\mu_\theta(s_t))_{\mathrm{steer}} - a^{\mathrm{sh}}_{t,\mathrm{steer}}\bigr)^2}
       {\sum_t \alpha_t},
  \qquad
  \mathcal{L} \;=\; \mathcal{L}_{\mathrm{PPO}} + \lambda_{\mathrm{BC}}\,\mathcal{L}_{\mathrm{BC}},
\end{equation}
where \(\mu_\theta\) is the actor mean, \(\tanh(\mu_\theta(s_t))_{\mathrm{steer}}\) the policy's
deterministic steering output, \(a^{\mathrm{sh}}_{t,\mathrm{steer}}\) the steering component of the
shield's executed action, and \(\alpha_t\in[0,1]\) the per-step shield mixing coefficient of
Equation~\eqref{eq:alpha-blend}, so that each step is cloned in proportion to how strongly the
shield intervened (\(\alpha_t = 0\) on unshielded steps contributes nothing). The total loss adds
this term to the masked PPO surrogate \(\mathcal{L}_{\mathrm{PPO}}\) of
Section~\ref{sec:design-agent-ppo} with weight \(\lambda_{\mathrm{BC}}\ge0\).
The resulting learning rule is clean: PPO on the on-policy steps, behavioural cloning of the
shield on the shield-corrected steps, and the critic on all steps. The coefficient
\(\lambda_{\mathrm{BC}}\) is annealed to zero so the shield teaches strongly early and the agent
is released to its own policy late, a curriculum analogous to learning from
intervention~\cite{ross2011dagger}. Critically, the cloning target in
Equation~\eqref{eq:bc} is restricted to the \emph{steering} dimension; the reason is analysed in
Section~\ref{sec:dep-collapse}.

\subsection{Result: the agent learns to drive unaided}

The scheme is applied as a warm start from the competent shielded policy. The effect on the
shield-off probe is dramatic (Figure~\ref{fig:probe}): off-road episodes, the dominant unaided
failure, fall from \(\approx90\%\) to \(\approx6\%\), and the shield-off success rate rises from
\(0\%\) to \(50\text{--}75\%\) within a few hundred episodes, where the baseline curve stays flat
at zero. The agent has learned the shield's lane-keeping skill.

\begin{figure}[H]
  \centering
  \includegraphics[width=0.95\textwidth]{imgs/shield_off_probe}
  \caption{Shield-off probe (the true learning curve). Without the teacher scheme (grey) the
           agent never learns to drive unaided: it leaves the road within \(\approx40\,\mathrm{m}\)
           on essentially every episode. With the steering behavioural-cloning teacher (green),
           shield-off success climbs to \(50\text{--}75\%\) and off-road collapses to a few per
           cent.}
  \label{fig:probe}
\end{figure}

A final shield ablation confirms the result on identical, seed-aligned scenarios with \(60\) NPCs
(Table~\ref{tab:dep-ablation} and Figure~\ref{fig:ablation}). The headline number is that the
\emph{shield-off} success rate (\(66\%\)) \emph{matches} the \emph{shield-on} rate (\(66\%\)):
the agent no longer needs the shield to function. What the shield still buys is safety: it
reduces the crash rate from \(16\%\) to \(12\%\) and the off-road rate from \(18\%\) to \(8\%\),
with the combined unsafe rate falling from \(34\%\) to \(20\%\), at the cost of \(12\%\)
stuck episodes induced by the shield's braking. The shield has gone from a crutch the agent
could not function without to a safety net it does not need in order to drive.

\begin{table}[H]
  \centering
  \caption{Shield ablation after weaning (\texttt{learn\_fix2}), \(50\) episodes,
           \(60\) NPCs, seed-aligned scenarios.}
  \label{tab:dep-ablation}
  \small
  \begin{tabular}{@{}lccccc@{}}
    \toprule
    \textbf{Configuration} & \textbf{Success} & \textbf{Crash} & \textbf{Off-road} &
    \textbf{Stuck} & \textbf{Shields/ep} \\
    \midrule
    Shield OFF (\texttt{none})      & \(66\%\) & \(16\%\) & \(18\%\) & \(0\%\) & \(0.0\) \\
    Shield ON (\texttt{adaptive})   & \(66\%\) & \(12\%\) & \(8\%\)  & \(12\%\) & \(65.6\) \\
    \bottomrule
  \end{tabular}
\end{table}

\begin{figure}[H]
  \centering
  \includegraphics[width=\textwidth]{imgs/shield_ablation}
  \caption{Ablation of the weaned agent (\texttt{learn\_fix2}, \(60\) NPCs). Shield-off and
           shield-on success are matched at \(66\%\); the shield's remaining value is reducing
           the combined crash\,+\,off-road rate from \(34\%\) to \(20\%\).}
  \label{fig:ablation}
\end{figure}

Figure~\ref{fig:weaned-comparison} contextualises this result against the pre-weaning
baseline (\texttt{FullSendC1\_p4}) across all three shield configurations under identical
conditions (20~NPCs, seed-aligned scenarios).

\begin{figure}[H]
  \centering
  \includegraphics[width=\textwidth]{imgs/weaned_vs_baseline}
  \caption{Baseline (\texttt{FullSendC1\_p4}, top row) versus weaned (\texttt{learn\_fix2},
           bottom row), 50 seed-aligned episodes, 20~NPCs. Left: outcome distribution per
           shield type. Centre: crash and off-road rates. Right: average distance and shield
           interventions per episode. The baseline agent requires the adaptive shield to
           drive at all (0\% success without it); the weaned agent reaches 74\% success
           with no shield and matches that rate under all three configurations.}
  \label{fig:weaned-comparison}
\end{figure}

\section{Failure analysis: why only the steering}\label{sec:dep-collapse}

Restricting the cloning target in Equation~\eqref{eq:bc} to the steering dimension is not a
detail but the difference between success and collapse. When the throttle/brake dimension is
included, the policy collapses to a stationary ``parking'' equilibrium
(Figure~\ref{fig:collapse}): the mean speed falls to near zero and the success rate to zero.

The mechanism is that \emph{the shield's throttle output is always an emergency brake}, as seen in
the lateral-recovery and emergency-brake commands of Section~\ref{sec:design-adaptive-shield}. Unlike
its steering, which demonstrates good lane keeping, the shield never demonstrates good
\emph{longitudinal} driving; it only ever demonstrates slamming the brakes. Cloning that signal
teaches the agent to brake, and because behavioural cloning is a \emph{direct} supervised
gradient on the policy mean it overpowers the reward-mediated incentives that would otherwise
keep the agent moving. A narrower variant that cloned the brake only on \emph{dynamic} (vehicle)
interventions was also tried; it too collapsed, as soon as the curriculum raised the traffic
density the now-frequent vehicle interventions taught systematic over-braking
(Figure~\ref{fig:collapse}). The lesson generalises: \textbf{the shield's steering is a good
teacher, but its throttle is never one}, because it encodes a last-resort reflex rather than a
driving policy.

\begin{figure}[H]
  \centering
  \includegraphics[width=0.85\textwidth]{imgs/throttle_collapse}
  \caption{Imitating the shield's throttle collapses the policy to parking. Cloning the steering
           only (green) is stable; cloning the throttle on every shielded step (orange) or only on
           dynamic interventions (red) drives the mean speed to zero, the latter precisely when
           the curriculum increases the traffic density.}
  \label{fig:collapse}
\end{figure}

\section{Discussion}\label{sec:dep-discussion}

The two approaches share a target but differ in what they treat as the cause of dependence. The
LLM meta-controller treats it as a \emph{scheduling} problem, asking how aggressively the
shield should filter, and the experiment of Section~\ref{sec:dep-llm} shows that this is the wrong
variable: loosening the filter neither teaches the agent nor reduces its reliance, it only
exposes it to more risk. The behavioural-cloning teacher treats dependence as a \emph{learning}
problem, since the agent never received gradient on the states the shield handled, and resolves
it by supplying that gradient. Equation~\eqref{eq:bc} is, in effect, the missing half of the
masked objective.

Two corollaries are worth recording. First, the safety governor of
Section~\ref{sec:dep-llm} is a sound and reusable component independently of the LLM: it never
permitted an unsafe run, and the ``proposer + deterministic veto'' pattern is a safe way to put a
learned controller in a safety-critical loop even when, as here, the proposer turns out to be
weak. Second, the residual gap after weaning is collision avoidance (the unsafe rate
remains at \(34\%\) without the shield against \(20\%\) with it); since the shield is a poor
longitudinal
teacher, closing that gap is better posed as a reward-shaping problem (a time-to-collision or
following-distance penalty that lets the agent learn its \emph{own} gentle avoidance) than as a
further imitation target, and is left as future work.

\clearpage
